\chapter{Itching Pruritus}

\section*{Definition}
Unpleasant sensation provokes desire to scratch may localized generalized acute chronic primary secondary cause diverse etiologies dermatological systemic neurological psychological vascular mechanical factors

\section*{Pathophysiology}
Neurotransmitter release substance P histamine cytokines nerve fiber activation itch-specific C fibers pruriceptors central processing spinal cord brain perception scratching provides temporary relief but damages skin creates itch-scratch cycle perpetuating condition

\section*{Causes}
\begin{itemize}
  \item Dry skin xerosis most common winter dry climates aging reduced sebum production
  \item Allergic contact dermatitis poison ivy nickel cosmetics chemicals urticaria hives angioedema
  \item Viral exanthems measles rose chickenpox HSV VZV
  \item Drug eruption Stevens Johnson DRESS syndrome serum sickness-like reaction
  \item Autoimmune cutaneous lupus vasculitis psoriasis pemphigus pemphigoid
\end{itemize}

\section*{When to See a Doctor}
See your doctor if:
\begin{itemize}
  \item Severe or worsening symptoms
  \item Symptoms lasting more than Systemic mucosal surfaces fever rapidly progressive blister mucosal erosions hospital rule out drug-induced life-threatening syndromes SJS TEN scalded skin syndrome widespread unexplained itching jaundice dark urine pale stool suggest cholestatic liver disease renal failure lymphoma thyroid dysfunction diabetes neuropathy psychiatric depression anxiety psychogenic pruritus
  \item Accompanied by other concerning signs
\end{itemize}

\section*{Related Symptoms}
\begin{itemize}
  \item Itching sore warmth scaling crust oozing blister rupture peeling thickening pigmentation pain bacterial superinfection fever joint eye involvement rash lesions bumps blisters
\end{itemize}
\textbf{Important:} Always consider serious underlying conditions when evaluating persistent itching pruritus

\section*{Self-Care / Home Management}
\begin{itemize}
  \item Keep the affected area clean and dry; wash gently with mild, fragrance-free soap and lukewarm water.
  \item Avoid scratching — use cool compresses, calamine lotion, or OTC hydrocortisone cream to relieve itching.
  \item Apply moisturizer regularly to maintain skin barrier integrity, especially after bathing.
  \item Identify and avoid potential triggers (irritants, allergens, excessive heat, tight clothing).
  \item Seek medical evaluation if the rash spreads rapidly, develops blisters, shows signs of infection (pus, increasing redness, warmth), or is accompanied by fever.
\end{itemize}

\section*{How Your Doctor Will Evaluate You}
\begin{itemize}
  \item History includes onset, location, progression, associated symptoms (itching, pain, fever), exposure to triggers, and prior skin conditions.
  \item Physical examination assesses morphology (macules, papules, vesicles, etc.), distribution, and extent of skin involvement.
  \item Diagnostic tests may include skin scraping for fungal examination, patch testing for allergic contact dermatitis, or biopsy for suspicious lesions.
  \item Laboratory studies (CBC, ESR, autoimmune serologies) are obtained when systemic disease is suspected.
\end{itemize}

\section*{Treatment Options}
\begin{itemize}
  \item Topical therapies (corticosteroids, antifungals, antibiotics, emollients) are first-line for most localized skin conditions.
  \item Systemic therapy (oral antihistamines, antibiotics, antifungals, or immunosuppressants) is reserved for extensive or refractory cases.
  \item Phototherapy or biologic agents may be indicated for chronic inflammatory skin diseases (psoriasis, eczema).
  \item Referral to dermatology is appropriate for uncertain diagnoses, treatment failure, or skin lesions suspicious for malignancy.
\end{itemize}

\clearpage
